\begin{proposition}[A target-only sufficient positivity hypothesis]
\label{c:threefold-target-positivity}
\status{Proved conditional refinement; the general primitive pairing remains unresolved}
Retain the strictly semistable threefold setting and fixed
$\ell\ne\operatorname{char}k$ of
Proposition~\ref{c:threefold-pairing-kernel}, and put
$V=Y^{(1)}$, $S=Y^{(2)}$.
After a finite extension of the finite residue field, we may suppose
that the geometric components are individually defined.
Assume the Hodge standard conjecture
$I(S\times S,L\boxplus L)$ and equality of $\ell$-homological
and numerical equivalence for rational cycles on $V\times S$.
Here the first assumption includes the products of distinct surface
components. Then the primitive restriction image
$R_{\rm prim}$ has nondegenerate intersection pairing.
Consequently weight--monodromy holds on $H^3$ of the generic
threefold. No Lefschetz standard conjecture on $V$ is required
in this sufficient package.

Alternatively, suppose that for each component $V_i$ an inverse to
\[
 L_i:H^2(V_i)\longrightarrow H^4(V_i)(1)
\]
is induced by an algebraic correspondence.
Choose smooth ample surfaces $j_i:T_i\hookrightarrow V_i$ with
classes $m_iL_i$, and put $T=\coprod_iT_i$.
Then the same conclusion follows from
$I(S\times S,L\boxplus L)$ and
homological equals numerical equivalence on $T\times S$.
Only the relevant homogeneous degree-two mixed trace pairings
are needed in either formulation; full equality of the two
equivalences on the displayed products is a sufficient hypothesis.
\end{proposition}

\begin{proof}
Rational algebraic degree projectors over a finite field are available
from Frobenius and
\citet[Theorem~2(1), pp.~76--77]{c_katzmessing1974}.
We therefore work with the homogeneous degree-two homological
objects $H_V,H_S$ and the algebraic restriction morphism
$f=\rho:H_V\to H_S$.
The primitive projector $p$ on the surface target is algebraic:
on each component it is the degree-two projector minus
\[
 x\longmapsto
 \frac{\operatorname{Tr}(xL)}{\operatorname{Tr}(L^2)}\,L.
\]

We first justify the target algebra and its precise pairing.
The Lefschetz standard conjecture for surfaces is known and is
stable under products. For the latter assertion, the normalized
lowering operator on a product is
$\Lambda_1\otimes1+1\otimes\Lambda_2$.
The degree operator and the $\mathfrak{sl}_2$ Casimir have finitely
many rational eigenvalues; polynomial interpolation recovers
each primitive-string level and hence the usual unnormalized
lowering operator. Thus the product lowering is algebraic.
It follows that $B(S\times S)$, and hence $A(S\times S)$, holds.
Combined with the assumed $I(S\times S)$, this gives $D(S\times S)$
as in \citet[\S4.1, pp.~14--15]{c_msaito2003positivity}.
Consequently
\[
 A=\End(H_S)
\]
identifies with the homogeneous corner of the numerical
correspondence algebra. It is finite-dimensional over $\Q$ and
semisimple by \citet[Theorem~1, pp.~448--449]
{c_jannsen1992semisimple}.
Finite-dimensionality over $\Q$ is obtained this way, not inferred
from the dimension of its $\Q_\ell$ realization.

On all of $H^2(S)$ we use the Lefschetz-modified pairing
$\Phi(x,y)=\operatorname{Tr}(x\cup *y)$, with Saito's signed
Lefschetz involution. Its adjoint $a\mapsto a^\dagger$ is algebraic,
and the assumed product positivity gives
\begin{equation}
\label{eq:c-target-positive-trace}
 \operatorname{Tr}(a^\dagger a)>0\qquad(0\ne a\in A),
\end{equation}
by \citet[equation~(4.1.3), p.~15]{c_msaito2003positivity}.
On primitive $H^2$ one has $*=-1$, so $\Phi$ is the negative
of the raw intersection pairing. They have the same radical
there. The raw intersection pairing on all of $H^2(S)$ is
not asserted to have positive adjoint trace.

For algebraic homological correspondences, put
\[
 M=\Hom(H_V,H_S),\qquad N=\Hom(H_S,H_V).
\]
The pairing
\[
 M\times N\longrightarrow\Q,\qquad
 (u,v)\longmapsto\operatorname{Tr}(uv)
\]
is nondegenerate in its first variable under $D(V\times S)$.
Indeed, in the covariant cohomological convention the two
classes have complementary codimensions $3$ and $2$ on
$V\times S$. Their composition trace is their intersection
number. Projecting a complementary class onto its homogeneous
degree-two block does not change its pairing with $u$.
Thus vanishing against every reverse block makes $u$
numerically zero, and the hypothesis makes its cohomological
class zero.

The subspace $J=fN\subset A$ is a right ideal.
Since $A$ is semisimple, $J=eA$ for an idempotent $e\in A$.
There is $g_0\in N$ with $e=fg_0$, so
$\operatorname{im}e\subset\operatorname{im}f$.
For every $g\in N$ one has $(1-e)fg=0$.
The mixed trace pairing therefore gives $(1-e)f=0$, whence
$\operatorname{im}f=\operatorname{im}e$.
This argument requires no semisimplicity of $\End(H_V)$.

The finite additive Karoubi category generated by $H_S$
is semisimple, and its faithful realization is exact because
its exact sequences split. The kernel of
\[
 (1-p)e:\operatorname{im}e\longrightarrow H_S
\]
therefore has an algebraic idempotent projector $r$ and
realizes exactly
$\operatorname{im}f\cap\operatorname{im}p=R_{\rm prim}$.
We have taken the actual intersection, not replaced it
by $p(\operatorname{im}f)$.
Its radical for $\Phi$ is
$\operatorname{im}r\cap\ker r^\dagger$, another algebraic
subobject. If nonzero, it has a nonzero idempotent projector $q$.
The pairing vanishes on $\operatorname{im}q$ with itself,
so $q^\dagger q=0$, contradicting
\eqref{eq:c-target-positive-trace}.
This proves the required nondegeneracy. The dimension-three
criterion of Proposition~\ref{c:threefold-pairing-kernel}
now proves the asserted weight--monodromy statement.

For the alternative package, the algebraic maps
\[
 r_i=m_i^{-1}L_i^{-1}j_{i*}:H^2(T_i)\longrightarrow H^2(V_i)
\]
satisfy $r_i j_i^*=1$ by the projection formula.
In this composition $j_{i*}$ has target $H^4(V_i)(1)$,
which is exactly the source of the supplied inverse.
Thus $r=\bigoplus_i r_i$ is surjective and
$\rho r$ has exactly the image of $\rho$.
Apply the preceding proof to the two surface objects $H_T,H_S$
using the assumed mixed pairing on $T\times S$.
A finite constant-field extension for the components or smooth
sections is harmless: the geometric spaces and maps are the same,
and the Frobenius eigenvalues are replaced by their powers,
which preserves all their weights.
\end{proof}

This sufficient package isolates an input that the known results
for an individual surface do not supply.
$I(S)$ controls divisor classes, whereas the required
positive adjoint trace involves codimension-two correspondence
classes on the fourfold $S\times S$.
Likewise $D(S)$ for divisors does not supply $D(T\times S)$.
A supplied degree-two inverse Hard Lefschetz on a threefold
reduces the mixed product to products of surfaces; it does not
prove either retained product hypothesis.
